\documentclass[11pt,a4paper]{article}

% --- Encoding & Language ---
\usepackage[utf8]{inputenc}
\usepackage[T1]{fontenc}
\usepackage[english]{babel}

% --- Page Geometry ---
\usepackage[margin=1in]{geometry}

% --- Typography ---
\usepackage{microtype}
\usepackage{tgtermes}
\renewcommand{\familydefault}{\rmdefault}

% --- Math ---
\usepackage{amsmath}
\usepackage{amssymb}

% --- Tables ---
\usepackage{booktabs}
\usepackage{tabularx}
\usepackage{longtable}
\usepackage{array}

% --- Code listings ---
\usepackage{listings}
\usepackage{xcolor}

\definecolor{codebg}{HTML}{F5F5F5}
\definecolor{codeframe}{HTML}{CCCCCC}
\definecolor{keyword}{HTML}{0000FF}
\definecolor{comment}{HTML}{008000}
\definecolor{string}{HTML}{A31515}

\lstset{
  backgroundcolor=\color{codebg},
  frame=single,
  rulecolor=\color{codeframe},
  basicstyle=\ttfamily\small,
  keywordstyle=\color{keyword}\bfseries,
  commentstyle=\color{comment}\itshape,
  stringstyle=\color{string},
  breaklines=true,
  breakatwhitespace=true,
  tabsize=4,
  showstringspaces=false,
  numbers=none,
  xleftmargin=4pt,
  xrightmargin=4pt,
  aboveskip=8pt,
  belowskip=8pt,
}

\lstdefinestyle{bash}{
  language=bash,
  morekeywords={python,SAM_CHECKPOINT},
}

\lstdefinestyle{python}{
  language=Python,
  morekeywords={self,None,True,False},
}

% --- Links ---
\usepackage[colorlinks=true,linkcolor=blue!60!black,citecolor=blue!60!black,urlcolor=blue!60!black]{hyperref}

% --- Graphics ---
\usepackage{graphicx}

% --- Lists ---
\usepackage{enumitem}
\setlist[itemize]{topsep=4pt, itemsep=2pt, parsep=0pt}
\setlist[enumerate]{topsep=4pt, itemsep=2pt, parsep=0pt}

% --- Header/Footer ---
\usepackage{fancyhdr}
\pagestyle{fancy}
\fancyhf{}
\fancyhead[L]{\small Technical Report}
\fancyhead[R]{\small Object-Aware Image Reconstruction Pipeline}
\fancyfoot[C]{\thepage}
\renewcommand{\headrulewidth}{0.4pt}

% --- Title ---
\title{\textbf{Technical Report: Implementation of the\\Object-Aware Image Reconstruction Pipeline}\\[12pt]
\large Accompanying: \textit{Object-Aware Image Reconstruction: Comparative Strategies\\for Selective Enhancement in Latent Diffusion Models}}

\author{
  Oliver Tong\\
  \small Horace Mann School, Bronx, NY\\
  \small \texttt{oliver\_tong@horacemann.org}
  \and
  Greg Shakhnarovich\\
  \small Toyota Technological Institute at Chicago\\
  \small \texttt{greg@ttic.edu}
}

\date{June 2026}

\begin{document}

\maketitle

\begin{abstract}
This technical report documents the software implementation of the object-aware image reconstruction framework. The codebase consists of seven Python scripts that implement four reconstruction methods (DOD-FADE, DOD-SAM, LOL-FADE, LOL-SAM), compute quantitative evaluation metrics, and generate ablation analyses for the stylization phenomenon. We describe each script's role, its algorithmic mechanics, inter-script dependencies, and provide instructions for reproducing all results reported in the accompanying paper.
\end{abstract}

\tableofcontents
\newpage

% ============================================================================
\section{Overview}

The repository implements the complete experimental framework for object-aware image reconstruction using latent diffusion models. The codebase is organized into seven scripts, each responsible for a distinct phase of the pipeline:

\begin{table}[h]
\centering
\small
\caption{Repository script inventory.}
\label{tab:scripts}
\begin{tabularx}{\textwidth}{lX}
\toprule
\textbf{Script} & \textbf{Purpose} \\
\midrule
\texttt{process\_image.py} & Core reconstruction pipeline (all 4 methods) \\
\texttt{evaluate\_image\_blending.py} & Per-image composite quality evaluation ($Q_{\text{img}}$) \\
\texttt{eval\_metrics.py} & Aggregate metrics (FID, LPIPS) computation \\
\texttt{compute\_iou\_direct.py} & IoU between method masks and GrabCut reference \\
\texttt{augment\_metrics\_with\_masks.py} & Boundary energy and background smoothness augmentation \\
\texttt{evaluate\_weight\_sensitivity.py} & Weight sensitivity robustness analysis \\
\texttt{generate\_stylization\_ablation.py} & Stylization ablation study \\
\bottomrule
\end{tabularx}
\end{table}

% ============================================================================
\section{Core Reconstruction Pipeline: \texttt{process\_image.py}}

\subsection{Architecture}

The core pipeline is implemented via two primary classes and a set of free functions.

\subsubsection{ImageProcessor Class}

Encapsulates model loading and inference:
\begin{itemize}
  \item \textbf{Object detection}: Faster R-CNN with ResNet-50-FPN backbone (torchvision pretrained weights), confidence threshold $\tau = 0.8$.
  \item \textbf{VAE}: Stable Diffusion v1.4 \texttt{AutoencoderKL} loaded from local \texttt{./models} directory via the \texttt{diffusers} library (uses safetensors format on Windows, pickle \texttt{.bin} on other platforms). Provides \texttt{encode()} and \texttt{decode()} methods that map between pixel space $\mathbb{R}^{H \times W \times 3}$ and latent space $\mathbb{R}^{4 \times h \times w}$ where $(h, w) = (H/8, W/8)$.
  \item \textbf{Device selection}: Automatic CUDA/MPS/CPU selection at initialization (supports Apple Silicon via MPS backend).
\end{itemize}

\subsubsection{ImageDisplayer Class}

Provides matplotlib-based visualization for side-by-side comparison of original and reconstructed images. Configurable grid layout via \texttt{rows} and \texttt{cols} parameters.

\subsubsection{Key Free Functions}

\begin{itemize}
  \item \texttt{get\_fade\_mask(width, height, fade\_px)}: Generates a heuristic edge-distance alpha mask in pixel space using:
  \begin{equation}
    \alpha(x,y) = \left[\text{clip}\left(\frac{d_{\text{edge}}(x,y)}{r_{\text{fade}}},\; 0,\; 1\right)\right]^{1.5}
  \end{equation}

  \item \texttt{get\_fade\_mask\_latent(width, height, fade\_px, device)}: PyTorch tensor version of the fade mask for latent-space operations.

  \item \texttt{get\_sam\_mask(image, box, sam\_checkpoint, sam\_model\_type, device)}: Runs SAM inference using the specified bounding box as a prompt. Auto-detects model type from checkpoint filename (\texttt{vit\_h} vs \texttt{vit\_b}). Device defaults to CUDA/MPS/CPU auto-detection.

  \item \texttt{tensor\_to\_pil(tensor)}: Converts a VAE output tensor (range $[-1, 1]$) to a PIL Image (range $[0, 255]$) via:
  \begin{equation}
    I_{\text{pixel}} = \frac{\text{clamp}(z, -1, 1) + 1}{2} \times 255
  \end{equation}
\end{itemize}

\subsection{Method Implementations}

\subsubsection{DOD-FADE: \texttt{process\_image\_dod\_fade()}}

Implements pixel-space reconstruction with fade-mask blending:

\begin{enumerate}
  \item \textbf{Background synthesis}: The input image is downscaled by factor 0.25, encoded into latent space, decoded back, then upscaled to original resolution. The encode-decode cycle induces smoothing.
  \item \textbf{Subject extraction}: Faster R-CNN detects objects; each ROI is cropped from the original image (original pixels, not re-encoded).
  \item \textbf{Compositing}: Each original subject crop is pasted onto the smoothed background using a fade mask (8\% of minimum ROI dimension as the fade radius, power exponent 1.5). PIL's \texttt{Image.paste()} performs the alpha compositing.
\end{enumerate}

\subsubsection{DOD-SAM: \texttt{process\_image\_dod\_sam()}}

Identical to DOD-FADE except the mask is generated by SAM rather than the heuristic fade function. The original subject pixels (not re-encoded) are pasted onto the smoothed background using the SAM mask. SAM is invoked with the bounding box as a prompt (\texttt{multimask\_output=False}). Falls back to fade mask if the \texttt{SAM\_CHECKPOINT} environment variable is unset or SAM inference fails.

\subsubsection{LOL-FADE and LOL-SAM: \texttt{process\_image\_lol()}}

Implements latent-space blending via the \texttt{mask\_type} parameter (\texttt{"fade"} or \texttt{"sam"}):

\begin{enumerate}
  \item \textbf{Latent background}: Input is downscaled by 0.25, encoded to latent, then the latent is upscaled back to full latent resolution using \texttt{F.interpolate} with \texttt{scale\_factor=1/scale} (bilinear mode):
  \begin{equation}
    z_{\text{bg,interp}} = \text{Upsample}(E(\text{Downscale}(I, 0.25)),\; \text{scale}=4)
  \end{equation}

  \item \textbf{Subject encoding}: Each detected ROI is cropped and encoded directly (without downscaling), producing a subject latent on the same grid:
  \begin{equation}
    z_j = E(\text{Crop}(I, B_j))
  \end{equation}

  \item \textbf{Latent blending}:
  \begin{itemize}
    \item \textbf{FADE mode}: Generates a fade mask at latent resolution and blends:
    \begin{equation}
      z_{\text{merged}} = z_{\text{bg}} \cdot (1 - M_{\text{fade}}) + z_j \cdot M_{\text{fade}}
    \end{equation}

    \item \textbf{SAM mode}: SAM produces a pixel-space mask, which is downsampled to latent resolution via bilinear interpolation, normalized to $[0, 1]$, then used for the same linear blend:
    \begin{equation}
      z_{\text{merged}} = z_{\text{bg}} \cdot (1 - M_{\text{SAM}}^{\text{lat}}) + z_j \cdot M_{\text{SAM}}^{\text{lat}}
    \end{equation}
  \end{itemize}

  \item \textbf{Decoding}: The merged latent is decoded to pixel space in a single pass: $I_{\text{final}} = D(z_{\text{merged}})$.
\end{enumerate}

Key implementation detail: coordinate mapping uses integer division by 8 (\texttt{y1 // 8}, \texttt{x1 // 8}) to convert pixel-space bounding box positions to latent-space positions.

\subsection{Stylized Subject Selection}

When \texttt{saturate\_factor} $\neq 1.0$ (stylization experiments), \texttt{process\_image\_lol()} applies single-subject selection logic:
\begin{itemize}
  \item For \texttt{cat3.png}: selects the 2nd detected object (to isolate the cat rather than a background detection).
  \item For all other images: selects only the 1st (highest-confidence) detected object.
\end{itemize}
This ensures the stylization effect is applied to a single primary subject rather than all detections.

\subsection{Subject Tiling}

The \texttt{find\_subjects()} method handles large objects that exceed configurable \texttt{max\_width}/\texttt{max\_height} thresholds (default 1000 pixels) by splitting them into tiles. For an object of width $w$ and height $h$:
\begin{align}
  h_{\text{splits}} &= \lceil w / \text{max\_width} \rceil \\
  v_{\text{splits}} &= \lceil h / \text{max\_height} \rceil
\end{align}
Each tile is processed independently, ensuring stable VAE encoding regardless of object size.

\subsection{Command-Line Interface}

\begin{lstlisting}[style=bash]
python process_image.py <method> <image_path> <output_name>
\end{lstlisting}

Where \texttt{<method>} is one of: \texttt{dod\_fade}, \texttt{dod\_sam}, \texttt{lol\_fade}, \texttt{lol\_sam}.

% ============================================================================
\section{Per-Image Composite Quality Evaluation: \texttt{evaluate\_image\_blending.py}}

\subsection{Metric Components}

This script implements the per-image $Q_{\text{img}}$ metric:
\begin{equation}
  Q_{\text{img}} = 0.45 \cdot E + 0.35 \cdot B + 0.20 \cdot O
\end{equation}

\paragraph{Edge Blending Score ($E$).}
Computes Sobel gradient magnitude across the boundary band (4-pixel morphological dilation minus erosion) and applies exponential decay:
\begin{equation}
  E = \exp\left(-\frac{\bar{g}_{\text{band}}}{50}\right)
\end{equation}
where $\bar{g}_{\text{band}}$ is the mean gradient within the boundary band. Score range $[0, 1]$; higher indicates smoother boundary transitions.

\paragraph{Background Smoothness Score ($B$).}
Computes the Laplacian of the grayscale image restricted to the background region:
\begin{equation}
  B = \exp\left(-\frac{\overline{|\nabla^2 I|}_{\text{bg}}}{30}\right)
\end{equation}
Higher values indicate fewer high-frequency artifacts in the background.

\paragraph{Object Preservation Score ($O$).}
Computes SSIM between the original and reconstructed image, restricted to the bounding box of the foreground mask. Uses \texttt{skimage.metrics.structural\_similarity}.

\subsection{Mask Extraction}

Rather than using the method's own mask, this script generates an independent reference mask using \textbf{GrabCut} (\texttt{cv2.GC\_INIT\_WITH\_RECT} mode). A 5\% border is treated as background prior, and 5 iterations are run. The binary mask is thresholded on \texttt{GC\_FGD | GC\_PR\_FGD}.

\subsection{Evaluation Loop}

The script evaluates all four methods across the 11 benchmark images (\texttt{car1-3}, \texttt{cat1-3}, \texttt{dog1-2}, \texttt{horse1-3}). Method results are read from \texttt{images/output/<method>/<image>\_0.25.png}. It prints per-image scores and identifies the winner for each image.

\subsection{Outputs}

\begin{itemize}
  \item \texttt{results/evaluate\_image\_blending\_result.json}: Per-image scores (edge blending, background smoothness, object preservation, $Q_{\text{img}}$ composite) for each method, plus per-image winners. This is the primary source for Table~5 (per-image $Q_{\text{img}}$) in the paper.
  \item \texttt{results/metrics.json} (merged): Adds \texttt{boundary\_energy} (boundary error $E$), \texttt{background\_laplacian}, and \texttt{object\_ssim} aggregates so that downstream scripts (notably \texttt{eval\_metrics.py}) can compute $Q_{\text{agg}}$ with real boundary values.
\end{itemize}

% ============================================================================
\section{Aggregate Metrics: \texttt{eval\_metrics.py}}

\subsection{FID Computation}

Evaluates the four primary methods plus a raw \texttt{output} baseline (VAE reconstruction without method-specific blending). Uses a ResNet-50 feature extractor (avgpool features, 2048-dimensional) rather than the standard Inception-v3 network. Feature vectors are collected across all images per method, and FID is computed as:
\begin{equation}
  \text{FID} = \|\mu_1 - \mu_2\|^2 + \text{Tr}\left(\Sigma_1 + \Sigma_2 - 2(\Sigma_1 \Sigma_2)^{1/2}\right)
\end{equation}
The covariance matrix square root is computed via \texttt{scipy.linalg.sqrtm}. With only 11 images, FID estimates are inherently high-variance due to rank-deficient covariance estimation.

\subsection{LPIPS Computation}

Preferentially uses the \texttt{lpips} package with VGG backbone when available (images resized to $256 \times 256$, normalized to $[-1, 1]$). Falls back to a VGG-16 perceptual distance (MSE in feature space at layer 16, images resized to $224 \times 224$) if the package is not installed.

\subsection{Outputs}

\begin{itemize}
  \item \texttt{results/metrics.json} (merged): Adds \texttt{fid}, \texttt{lpips}, \texttt{composite} ($Q_{\text{agg}}$), and \texttt{per\_image\_lpips} fields. Preserves existing fields (e.g., \texttt{boundary\_energy} from \texttt{evaluate\_image\_blending.py}, \texttt{iou} from \texttt{compute\_iou\_direct.py}).
  \item \texttt{results/per\_image\_metrics.csv}: Per-image LPIPS values in flat CSV format (one row per image$\times$method pair).
  \item \texttt{results/quant\_table.tex}: \LaTeX{} snippet for direct inclusion in the paper.
\end{itemize}

\subsection{Distinction Between Per-Image Output Files}

Two files contain per-image results but they measure different things and come from different scripts:

\begin{table}[h]
\centering
\small
\begin{tabularx}{\textwidth}{lXX}
\toprule
 & \texttt{evaluate\_image\_blending\_result.json} & \texttt{per\_image\_metrics.csv} \\
\midrule
\textbf{Produced by} & \texttt{evaluate\_image\_blending.py} & \texttt{eval\_metrics.py} \\
\textbf{Metrics} & Edge blending, background smoothness, object preservation, $Q_{\text{img}}$ & LPIPS only \\
\textbf{Purpose} & Artistic composite quality (which method wins each image) & Pixel-level perceptual fidelity vs.\ original \\
\textbf{Format} & JSON with nested scores per method per image & Flat CSV (image, method, lpips) \\
\bottomrule
\end{tabularx}
\end{table}

These files are complementary with no overlap: one captures the three-component composite quality used to determine per-image winners (Table~5 in the paper), the other captures perceptual distance used in the aggregate comparison (Table~3).

% ============================================================================
\section{IoU Computation: \texttt{compute\_iou\_direct.py}}

\subsection{Mask Generation}

For each image, the script:
\begin{enumerate}
  \item Runs Faster R-CNN (on CPU) to obtain the top-confidence bounding box.
  \item Generates a \textbf{SAM mask} by running SAM ViT-H inference (hardcoded) with the bounding box as prompt.
  \item Generates a \textbf{FADE mask} using the same heuristic as \texttt{process\_image.py} and binarizes it at the 127 threshold.
\end{enumerate}

\subsection{IoU Calculation}

Both masks are placed into full-image-resolution arrays using \texttt{paste\_roi\_mask()}, which handles coordinate clamping and size mismatches. IoU is computed against ground-truth masks stored in \texttt{images/masks/}:
\begin{equation}
  \text{IoU} = \frac{|A \cap B|}{|A \cup B|}
\end{equation}
The same IoU value is reported for DOD-SAM and LOL-SAM (identical mask), and for DOD-FADE and LOL-FADE (identical mask).

% ============================================================================
\section{Boundary Error Augmentation: \texttt{augment\_metrics\_with\_masks.py}}

This script augments the \texttt{metrics.json} file with three additional per-method metrics. It reads method results from \texttt{images\_eval/<method>/} and ground-truth masks from \texttt{images/masks/}:

\begin{itemize}
  \item \textbf{Boundary Error ($E$)}: RMSE of image gradients within a 4-pixel morphological boundary band (8-pixel total width). Reported as raw value (not exponentially decayed); lower values indicate smoother seam integration. Note: the code variable is named \texttt{boundary\_energy} but corresponds to Boundary error $E$ in the paper.
  \item \textbf{Background Laplacian}: Mean absolute Laplacian in the background region (raw high-frequency energy).
  \item \textbf{Object SSIM}: Structural similarity restricted to the object bounding box.
\end{itemize}

It also regenerates the \LaTeX{} table snippet (\texttt{results/quant\_table.tex}) with the updated values.

% ============================================================================
\section{Weight Sensitivity Analysis: \texttt{evaluate\_weight\_sensitivity.py}}

\subsection{Design}

Tests robustness of the LOL-SAM $>$ LOL-FADE conclusion under three weighting schemes:

\begin{table}[h]
\centering
\small
\caption{Weighting schemes for sensitivity analysis.}
\begin{tabular}{lccc}
\toprule
\textbf{Scheme} & \textbf{Edge} & \textbf{Smoothness} & \textbf{Object Preservation} \\
\midrule
Original & 0.45 & 0.35 & 0.20 \\
Equal & 0.333 & 0.333 & 0.333 \\
Edge-only & 1.0 & 0.0 & 0.0 \\
\bottomrule
\end{tabular}
\end{table}

\subsection{Implementation}

Imports the metric functions directly from \texttt{evaluate\_image\_blending.py} and reads method results from \texttt{images\_eval/<method>/}. Re-computes composite scores under each weighting scheme. For each scheme, reports:
\begin{itemize}
  \item Per-image scores and winners.
  \item Method averages and ranking.
  \item Per-image win counts (out of 11).
\end{itemize}

\subsection{Key Result}

LOL-SAM wins 9/11 images under all three schemes, confirming that the preference is not an artifact of the weight distribution. The advantage over LOL-FADE actually strengthens under edge-only evaluation (mean difference widens from 0.052 to 0.088).

% ============================================================================
\section{Stylization Ablation Study: \texttt{generate\_stylization\_ablation.py}}

\subsection{Purpose}

Validates the paper's third contribution: that mask-alpha extrapolation beyond $[0, 1]$ produces controlled artistic stylization via VAE decoder saturation.

\subsection{Ablation Components}

\paragraph{Mask-scale sweep.}
Runs LOL-SAM at $s_\alpha \in \{0.0, 0.5, 1.0, 1.5, 2.0\}$. The key modification:
\begin{lstlisting}[style=python]
mask_tensor = mask_tensor / 255.0 * alpha  # alpha > 1 extrapolates
\end{lstlisting}
This pushes latent activations beyond the training distribution of the VAE decoder.

\paragraph{Latent statistics.}
For each $s_\alpha$, records the distribution of latent values immediately before decoding: mean, standard deviation, skewness, min/max, fraction outside $|z| > 1$, and fraction outside $|z| > 3$.

\paragraph{Pixel-space $\alpha > 1$ control.}
Multiplies RGB values by $s_\alpha$ directly (with clamping to $[0, 255]$), composited with a SAM mask. Tests whether simple pixel scaling reproduces the latent stylization effect.

\paragraph{Filter baselines.}
Applies unsharp masking, bilateral filtering, and CLAHE to the entire image for comparison.

\subsection{Stylization Strength Metric}

Each output variant is evaluated using a stylization strength metric:
\begin{equation}
  \text{Stylization} = \frac{\overline{|\nabla^2 I|}_{\text{fg, result}}}{\overline{|\nabla^2 I|}_{\text{fg, original}}}
\end{equation}
where $\overline{|\nabla^2 I|}$ is the mean absolute Laplacian (high-frequency energy) restricted to the foreground region. Values $> 1$ indicate increased texture relative to the original.

\subsection{Outputs}

\begin{itemize}
  \item 110 ablation images in \texttt{images/output/ablation/}.
  \item Full per-image results in \texttt{results/stylization\_ablation\_results.json}.
  \item Aggregated CSV summary in \texttt{results/stylization\_ablation\_summary.csv}.
\end{itemize}

% ============================================================================
\section{Dependencies and Environment}

\subsection{Python Dependencies}

\begin{table}[h]
\centering
\small
\caption{Required Python packages.}
\begin{tabularx}{\textwidth}{llX}
\toprule
\textbf{Package} & \textbf{Version} & \textbf{Purpose} \\
\midrule
\texttt{torch} & $\geq$ 1.13 & Tensor computation, model inference \\
\texttt{torchvision} & $\geq$ 0.14 & Faster R-CNN, transforms \\
\texttt{diffusers} & 0.4.2 & Stable Diffusion VAE (AutoencoderKL) \\
\texttt{segment-anything} & $\geq$ 1.0 & SAM segmentation \\
\texttt{Pillow} & $\geq$ 9.0 & Image I/O \\
\texttt{numpy} & $\geq$ 1.21 & Array operations \\
\texttt{opencv-python} & $\geq$ 4.5 & GrabCut, gradient computation \\
\texttt{scikit-image} & $\geq$ 0.19 & SSIM metric \\
\texttt{scipy} & $\geq$ 1.7 & FID matrix square root \\
\texttt{matplotlib} & $\geq$ 3.5 & Visualization \\
\texttt{lpips} & $\geq$ 0.1 & Perceptual similarity (optional) \\
\bottomrule
\end{tabularx}
\end{table}

\subsection{Model Weights}

\begin{table}[h]
\centering
\small
\caption{Pretrained model weights.}
\begin{tabularx}{\textwidth}{lXl}
\toprule
\textbf{Model} & \textbf{File} & \textbf{License} \\
\midrule
SD VAE v1.4 & \texttt{models/diffusion\_pytorch\_model.bin} & CreativeML Open RAIL-M \\
SAM ViT-H & \texttt{models/sam\_vit\_h\_4b8939.pth} & Apache 2.0 \\
SAM ViT-B & \texttt{models/sam\_vit\_b\_01ec64.pth} & Apache 2.0 \\
Faster R-CNN & torchvision hub (auto-downloaded) & BSD 3-Clause \\
\bottomrule
\end{tabularx}
\end{table}

Model files are tracked via Git LFS. Run \texttt{git lfs pull} after cloning.

\subsection{Environment Variables}

\begin{table}[h]
\centering
\small
\begin{tabularx}{\textwidth}{llX}
\toprule
\textbf{Variable} & \textbf{Default} & \textbf{Purpose} \\
\midrule
\texttt{SAM\_CHECKPOINT} & \texttt{models/sam\_vit\_b\_01ec64.pth} & Path to SAM model weights (ViT-B included; set to \texttt{sam\_vit\_h\_4b8939.pth} for ViT-H if downloaded) \\
\bottomrule
\end{tabularx}
\end{table}

% ============================================================================
\section{Reproduction Instructions}

\subsection{Default Parameters}

Unless noted otherwise, all reconstruction comparisons in the paper use:
\begin{itemize}
  \item \texttt{--scale 0.25} --- background downscaling factor (the VAE encodes a $4\times$ downscaled image, then the latent is upscaled back to full resolution).
  \item \texttt{--saturate\_factor 1.0} --- no mask-alpha extrapolation (standard reconstruction). Values $> 1.0$ are used only in the stylization experiments (Section~8).
\end{itemize}
Output images from these defaults are stored with the \texttt{\_0.25} suffix, e.g.\ \texttt{images/output/dod\_fade/car1\_0.25.png}.

\subsection{Generate Reconstruction Outputs (All 4 Methods)}

\begin{lstlisting}[style=bash]
cd scripts/

# Run all 4 methods for a single image (repeat for all 11 benchmark images)
for method in dod_fade dod_sam lol_fade lol_sam; do
    python process_image.py \
        --method $method \
        --input_img_path images/input/car1.png \
        --output_img_path "$method/car1_0.25" \
        --scale 0.25 \
        --max_subjects 1
done
\end{lstlisting}
Outputs are saved to \texttt{images/output/<method>/<image>\_0.25.png}.

\subsection{Run Evaluation Scripts}

\textbf{Important: scripts must be run in the order below.} \texttt{evaluate\_image\_blending.py} writes boundary energy, background smoothness, and object SSIM into \texttt{results/metrics.json}. \texttt{eval\_metrics.py} reads those boundary values to compute the correct $Q_{\text{agg}}$ composite score. Running \texttt{eval\_metrics.py} before \texttt{evaluate\_image\_blending.py} will produce $Q_{\text{agg}}$ values computed with a placeholder boundary term.

\begin{lstlisting}[style=bash]
cd scripts/

# Step 1: Per-image composite quality (Q_img) and boundary energy
python evaluate_image_blending.py

# Step 2: Aggregate metrics (FID, LPIPS) and composite Q_agg
python eval_metrics.py

# Step 3: IoU between method masks and GrabCut reference
python compute_iou_direct.py
\end{lstlisting}

\subsection{Weight Sensitivity Analysis}

\begin{lstlisting}[style=bash]
python evaluate_weight_sensitivity.py
\end{lstlisting}

\subsection{Stylization Ablation}

\begin{lstlisting}[style=bash]
SAM_CHECKPOINT=models/sam_vit_h_4b8939.pth python generate_stylization_ablation.py
\end{lstlisting}

% ============================================================================
\section{Design Decisions and Known Limitations}

\subsection{Design Decisions}

\begin{enumerate}
  \item \textbf{Module-level model instantiation}: \texttt{process\_image.py} instantiates \texttt{ImageProcessor} and \texttt{ImageDisplayer} at import time (lines 519--520). This was chosen for interactive notebook usage but causes side effects when importing the module. The ablation script (\texttt{generate\_stylization\_ablation.py}) reimplements the classes to avoid this.

  \item \textbf{Coordinate mapping via integer division}: The pixel-to-latent mapping \texttt{y1 // 8} introduces up to 7 pixels of alignment error. This is acceptable given the VAE's $8\times$ downsampling, which already blurs spatial boundaries.

  \item \textbf{SAM fallback to FADE}: The DOD-SAM method gracefully degrades to fade-mask blending if SAM is unavailable, ensuring the pipeline remains functional without GPU-intensive segmentation.

  \item \textbf{GrabCut as reference}: The evaluation uses GrabCut rather than manually annotated masks as the segmentation reference. This keeps evaluation fully automatic but means IoU measures inter-method agreement rather than absolute accuracy.
\end{enumerate}

\subsection{Known Limitations}

\begin{enumerate}
  \item \textbf{No batch processing}: Each image is processed independently; models are not shared across invocations in CLI mode (though they persist in module-level instances for scripted use).

  \item \textbf{Fixed hyperparameters}: The fade radius (8\%), power exponent (1.5), and detection threshold (0.8) are hardcoded. The paper justifies these choices empirically but does not expose them as CLI arguments.

  \item \textbf{Single-object priority in IoU}: \texttt{compute\_iou\_direct.py} uses only the highest-confidence detection for IoU computation, while \texttt{process\_image.py} processes all detections above threshold.

  \item \textbf{Memory}: SAM ViT-H requires approximately 2.5\,GB of VRAM. The pipeline will OOM on GPUs with less than 4\,GB when processing multiple subjects simultaneously.
\end{enumerate}

% ============================================================================
\section{Script Dependency Graph}

The following diagram shows the inter-script dependencies:

\begin{verbatim}
process_image.py  (standalone -- core pipeline, generates images/output/)
     |
     v
evaluate_image_blending.py  (reads images/output/*_0.25.png)
     |                       (writes evaluate_image_blending_result.json)
     |                       (merges boundary_energy into metrics.json)  <- MUST RUN FIRST
     v
eval_metrics.py  (reads boundary_energy from metrics.json)
     |            (computes FID/LPIPS, writes Q_agg with real boundary)
     v
compute_iou_direct.py  (standalone -- merges IoU into metrics.json)

evaluate_weight_sensitivity.py  (imports from evaluate_image_blending)

generate_stylization_ablation.py  (reimplements pipeline classes)
\end{verbatim}

\textbf{Execution order matters:} \texttt{evaluate\_image\_blending.py} must run before \texttt{eval\_metrics.py} so that boundary energy values are available for the $Q_{\text{agg}}$ computation. If \texttt{eval\_metrics.py} runs first, it will use a placeholder boundary value of 0.5 and print a warning.

% ============================================================================
\section{Correspondence Between Code and Paper Equations}

Table~\ref{tab:correspondence} maps each equation in the paper to its implementation location.

\begin{table}[h]
\centering
\small
\caption{Paper equation to code correspondence.}
\label{tab:correspondence}
\begin{tabularx}{\textwidth}{lX}
\toprule
\textbf{Paper Equation} & \textbf{Code Location} \\
\midrule
Eq.~1: $z = E(I)$ & \texttt{ImageProcessor.encode()} (line 82) \\
Eq.~2: $I' = D(z)$ & \texttt{ImageProcessor.decode()} (line 93) \\
Eq.~5: $d_{\text{edge}}$ & \texttt{get\_fade\_mask()} (lines 140--154) \\
Eq.~6: $M_{\text{FADE}}$ & \texttt{get\_fade\_mask()} (lines 152--157) \\
Eq.~7: Alpha compositing & \texttt{process\_image\_dod\_fade()} (line 235, \texttt{PIL.paste}) \\
Eq.~8: SAM mask generation & \texttt{get\_sam\_mask()} (lines 176--203) \\
Eqs.~10--12: Latent background & \texttt{process\_image\_lol()} (lines 353--360) \\
Eq.~14: Latent masking & \texttt{paste\_latend\_with\_fade()} (lines 294--306) \\
Eq.~15: Latent integration & \texttt{paste\_latend\_with\_sam()} (lines 308--332) \\
Eq.~16: Final decode & \texttt{process\_image\_lol()} (line 393) \\
$Q_{\text{img}}$ composite & \texttt{evaluate\_method()} in \texttt{evaluate\_image\_blending.py} (lines 167--171) \\
\bottomrule
\end{tabularx}
\end{table}

% ============================================================================
\section{Summary}

This codebase implements a complete experimental framework for object-aware image reconstruction using latent diffusion models. The modular design separates reconstruction (generation) from evaluation (analysis), enabling independent verification of reported results. All reconstruction comparisons use \texttt{scale=0.25} and \texttt{saturate\_factor=1.0} unless noted otherwise. All quantitative claims in the paper---including per-image composite quality scores, aggregate metrics, weight sensitivity analysis, and stylization ablation results---can be reproduced by running the scripts in the order described in Section~10: \texttt{evaluate\_image\_blending.py} must run before \texttt{eval\_metrics.py} to ensure correct $Q_{\text{agg}}$ computation with real boundary energy values.

\end{document}
